\subsection{F.1. Manuscript wording and statement map}
We quote the supplied AI transcription witnesses, not the absent manuscript
scans.  We omit struck-out words and retain inserted words where the resulting
reading is unambiguous.  Cote and scan numbers, rather than the manuscript's
own page numbers, follow the prompt's convention.  A quotation records the
question being addressed; it does not itself certify the mathematical claim.

\par\smallskip\noindent\textbf{A.1, cote 4:4.}
``Nous admettrons que ces relations restent valables dans le cas de la
cohomologie cristalline.''  The assertions are integral injectivity, torsion
identification and torsion-freeness of $E(Y)$.
Theorems~\ref{thm:A-high-degree} and \ref{thm:A-all-degree-ranges}
give the explicit sufficient degree range and two all-degree classes.
Proposition~\ref{prop:A-transpose-u2} types the missing (8.2.1): it is the
torsion Gysin transpose, and after the justified torsion identification its
composite is cup product by $d\,c_1(\mathcal O_X(1))$.
The full small-degree problem is retained in
Question~\ref{open:A-small-degree}.

\par\smallskip\noindent\textbf{A.2, cote 4:8, and the sheet 4:2.}
``Je pense que l'homomorphisme $v$ de (8.5.2) est injectif'' is settled by
Lemma~\ref{lem:A-R1} without a degree bound.
Theorems~\ref{thm:A-exact85} and \ref{thm:A-unit-count} then establish
the stable exact sequence, the dimension formulas and the unit-root count
for every smooth section.  The remark ``Sans doute (*) et (**) sont
transposés l'un de l'autre'' is made precise in
Proposition~\ref{prop:A-star-dual}; the adjoint operators are $F$ and $V$.

\par\smallskip\noindent\textbf{A.3, cote 4:9.}
The retained marginal assertion is ``Or cela est établi au n° 7'', beside
the eventual positivity of $\dim E(Y,\mathcal O_Y)_{\mathrm{ss}}$.
Theorem~\ref{thm:all-X-explicit-positive-rank} supplies an explicit
threshold, and Corollary~\ref{cor:eventual-hodge-parenthesis} settles
the eventual generic Hodge-level parenthesis.  The relative-ordinarity
prediction is refuted by the characteristic-two surface constructed above.
Proposition~\ref{prop:effective-generic-rank} computes the exact generic
rank for specified finite-field input; a structural rank formula for
arbitrary ambient $X$ remains an explicitly retained obligation.

\par\smallskip\noindent\textbf{A.4, cote 4:12.}
The question asks whether the image ``est-elle égale à
$\operatorname{Aut}(V)$ tout entier ?'', with an inserted question on
irreducibility.  Theorems~\ref{thm:pn-odd-full-groups} and
\ref{thm:pn-char2-all-groups} determine every projective-space case.
The latter refutes the unrestricted full-linear prediction in infinitely
many degrees.  Proposition~\ref{prop:exact-general-determinant-image}
determines the geometric and arithmetic determinant images for general
$X$ from the stable Hasse invariant.  The complete nonabelian image for
arbitrary $X$ remains open here.

\par\smallskip\noindent\textbf{A.5, cote 4:12.}
The requested consequence is that ``$\operatorname{card}Y(k')$ soit
congru mod $p$ à un élément de $\underline F_p$ donné à l'avance''.
Theorem~\ref{thm:all-X-residue-density} proves this for every ambient
variety in the question, with its explicit $d_{\rm res}$, in every
nonempty open.  Corollary~\ref{cor:pn-all-p-sharp-residue-threshold}
gives the sharp bound $d=n+2$ in projective space, including $d=4$
for planes.  The low-degree residue cases appear in the plane and
projective-space classifications.

\par\smallskip\noindent\textbf{A.6, cote 4:10--11.}
The expected comparison concerns the ``torsion de cohomologie
$p$-adique de la variété relevée''.
Theorem~\ref{thm:A-special-torsion} gives the exact special-fibre
identification and specialization obstruction.
Theorem~\ref{thm:A-BMS-counterexample} refutes the unrestricted
expectation; Corollary~\ref{cor:A-lifting-small-range} gives the
valid small-ramification implication.  The final ample-section search
in the counterexample is specified by a terminating construction,
not represented as an executed numerical search.
